%! Principles of Experimental Design — Unit 1, Lesson 1.6 — Student Packet Cover Sheet
% PILOT SHAPE (2026-09-06): modeled on AP Statistics lesson 1.4 minus its AP Practice page.
% THREE packet rows — Warm-Up, Guided Notes & Practice, Homework. The homework is the individual
% practice, started in class, and it is SCORED, so its cell is a \blank{}, never NA.
\documentclass[12pt]{article}
\usepackage{saar-article}
\usepackage{saar-boxes}
\usepackage{ltablex}
\keepXColumns
% Measured banner: the band grows to fit the title block, so a title long enough to wrap grows
% the band rather than running out of it (the stats course's \coverbanner; saar-article does not
% carry one, so it is defined here — every cover copies this block verbatim).
\newsavebox{\bannerbox}
\newlength{\bannerht}
\newlength{\coverbannerpad}
\setlength{\coverbannerpad}{0.16in}       % breathing room above and below the text
\newcommand{\coverbanner}[2]{%
  \sbox{\bannerbox}{%
    \begin{minipage}{\dimexpr\paperwidth-1.4in\relax}%
      \centering
      {\color{white}\textbf{\LARGE \CourseName}}\\[4pt]
      {\color{frostmid}\large\textbf{#1}}\\[6pt]
      {\color{white}\textbf{\large #2}}%
    \end{minipage}}%
  \setlength{\bannerht}{\dimexpr\ht\bannerbox+\dp\bannerbox+2\coverbannerpad\relax}%
  \noindent\begin{tikzpicture}[remember picture, overlay]
    \fill[cerulean] (current page.north west)
      rectangle ([yshift=-\bannerht]current page.north east);
    \node[anchor=north, inner sep=0pt]
      at ([yshift=-\coverbannerpad]current page.north) {\usebox{\bannerbox}};
  \end{tikzpicture}\par
  % The text body starts 0.6in below the page edge (geometry top=0.6in), so reserve only what
  % the band overruns that, plus a gap under the band.
  \vspace*{\dimexpr\bannerht-0.6in+0.16in\relax}%
}

\begin{document}

% Full-bleed cerulean banner, sized to the title block.
\coverbanner{Unit 1}{Lesson 1.6 \quad Principles of Experimental Design}

\vspace{0.06in}
% NAMESTRIP: this is the ONLY component that carries the name row.
\namedateperiod
\vspace{0.18in}

% GRADUAL RELEASE: the vocabulary is taught in the guided notes, so the targets name it
% plainly. The standard codes (PS.DC.3a-b, AFDA.DA.2c) stay in the teacher-facing plan.
\begin{learningtargetbox}
\small
\begin{enumerate}[leftmargin=1.5em, itemsep=5pt]
  \item \ldots{} tell an \textbf{experiment} from an observational study by asking who
        built the groups, and name the \textbf{experimental units}, the
        \textbf{treatment}, and the \textbf{control group}.
  \item \ldots{} check a design against the four principles ---
        \textbf{comparison, randomization, replication, control} --- and name the
        principle a changed plan breaks.
  \item \ldots{} explain the \textbf{placebo effect}, say who is \textbf{blinded} in a
        design, and split the units of a \textbf{randomized block design} into its
        blocks.
  \item \ldots{} justify why \textbf{random assignment} lets an experiment say
        \emph{caused} --- and why the size of the gap is not what earns that word.
\end{enumerate}
\end{learningtargetbox}

\vspace{0.14in}

\begin{tocbox}
\small
\renewcommand{\arraystretch}{1.5}
\begin{tabularx}{\linewidth}{@{} c l X r @{}}
\rowcolor{frost}
\textbf{\#} & \textbf{Component} & \textbf{Description} & \textbf{Score} \\
\midrule
1 & Warm-Up & Riverbend's Study Center experiment: the two group rates, the gap, and
             last quarter's gap the other way & \blank{1.2cm} \\
2 & Guided Notes \& Practice & Assigning instead of watching; the four principles; blinding
             and blocking; what buys the word \emph{caused} --- then the Harbor Point
             swim-time experiment worked together & \blank{1.2cm} \\
% Row 3 is the lesson's GRADED practice, started in class. Packet pages are the default; a
% Desmos activity may replace them on a given day, decided at Close & Assign. The row and its
% score cell never change.
3 & Homework & Millbrook Public Library runs the experiment this time: name the parts, check
             the principles, block on age, and explain why a smaller gap is the stronger
             evidence --- \textbf{scored, started in class, due the first class after two
             study halls} & \blank{1.2cm} \\
\midrule
  & \multicolumn{2}{r}{\textbf{Total}} & \blank{1.2cm} \\
\end{tabularx}
\end{tocbox}

\vspace{0.14in}

% Keep in Mind carries the lesson's CONTENT — the rule, the test, the trap — never its process.
\begin{remindbox}
\small
\textbf{In an experiment you assign --- you do not just watch.} Chance puts the
\textbf{experimental units} into groups and the researcher imposes a \textbf{treatment}
on at least one of them, while a \textbf{control group} shows what happens without it.
Check four principles: \textbf{comparison} (at least two groups), \textbf{randomization}
(chance picks, not people), \textbf{replication} (many units per group), and
\textbf{control} (everything else the same, so no \textbf{confounding variable} sneaks in).
Because chance built the groups, this is the one kind of study that earns the word
\emph{caused}. \textbf{Watch out:} a bigger gap is not stronger evidence. A gap between
groups that chose themselves is inflated by whoever chose; a smaller gap between groups
chance built is the treatment by itself.
\end{remindbox}

\end{document}
